\documentclass[../main.tex]{subfiles}

\begin{document}

\ifSubfilesClassLoaded{

    %\tableofcontents
    
    %\setcounter{chapter}{}
}{}

\chapter{ Mideum Test 2024 }
% 代靖涵 25120222201319
\begin{exercise}{}{}
袋子中有 $m$ 枚正品硬币和 $n$ 枚次品硬币 (次品硬币的两面都有国徽), 在袋中随机选取一枚硬币投掷 $r$ 次, 已知每次都得到国徽. 问这枚硬币是正品的概率.
\end{exercise}
\begin{proof}
    设 \(  A  \)为硬币是正品的事件,  设 \(  B  \)为投掷 \(  r  \)次每次都得到国徽的事件. 则题述事件为 \(  A|B  \), 概率为 \[
    \begin{aligned}
    P\left( A|B \right)&= \frac{P\left( AB \right)  }{P\left( B \right)  }\\ 
     &= \frac{P\left( B|A \right)P\left( A \right)   }{P\left( B|A \right)P\left( A \right) + P\left( B|A^{c} \right)P\left( A^{c} \right)    }    
    \end{aligned}
    \] 其中 \[
    P\left( A \right)= \frac{m }{m+ n },\quad P\left( A^{c} \right)= \frac{n }{m+ n }    
    \] \[
    P\left( B|A \right) = \frac{1 }{2^{r} },\quad P\left( B|A^{c} \right)= 1  
    \]于是 \[
    P\left( A|B \right)= \frac{\frac{1 }{2^{r} }\frac{m }{m+ n }   }{ \frac{1 }{2^{r} }\frac{m }{m+ n }  + \frac{n }{m+ n } }= \frac{m }{m+ 2^{r}n }   
    \]为此时这枚硬币是正品的概率.
\end{proof}
\begin{exercise}{}{}
一个靶子是半径为 $2$ 米的圆盘, 设击中靶上任一同心圆盘的点的概率与该圆盘的面积成正比, 并且每次射击都能击中靶, 以 $X$ 表示弹着点与圆心的距离, 试求随机变量 $X$ 的分布函数与密度函数.
\end{exercise}
\begin{proof}
     \[
     P\left( X\le r \right)=  \frac{ \pi r^{2} }{ 4 \pi }= \frac{r^{2} }{4 }  ,\quad  0\le r\le 2
     \] \[
     P\left( X\le r \right)= 0,\quad r< 0,\quad P\left( X\le r \right)= 1,\quad r> 2  
     \]于是 \(  X  \)的分布函数 \(  F_{X}  \)为 \[
     F_{X}\left( x \right)= \begin{cases} 0,&r< 0\\ 
      \frac{x^{2} }{4 },& 0\le r\le 2\\ 
       1,&r> 2  \end{cases}  
     \]  \(  X  \)的一个密度函数\(  p_{X}  \) 为 \(  F_{X}  \)的一个分段导数:  \[
     p_{X}\left( x \right)= \begin{cases} 0,&r< 0\\ 
      \frac{x }{2 },& 0\le r\le 2\\ 
       0,&r> 2  \end{cases}  
     \] 
\end{proof}
\begin{exercise}{}{}
设 $X \sim P(\lambda)$, 求证:
\[
F(n) = P(X \le n) = 1 - \int_{0}^{\lambda} \frac{x^n e^{-x}}{n!} dx
\]
\end{exercise}
\begin{proof}
     \[
     P\left( X= k \right)= \frac{ \lambda ^{k} }{k! }e^{- \lambda }  
     \] \[
     F\left( n \right)= \sum _{k= 0}^{n}\frac{ \lambda ^{k} }{k! }e^{- \lambda }  
     \]对于任意的 \(  n\ge 1  \), 我们有  \[
   \begin{aligned}
     \int_{0}^{ \lambda }\frac{x^{n}e^{-x} }{n! }\,\mathrm{d} x&= - \int_{0}^{ \lambda } \frac{x^{n} }{n! }\,\mathrm{d} e^{-x}  \\ 
      &=  -\left[ \frac{x^{n} }{n! }e^{-x} \right]_{0}^{ \lambda }+  \int_{0}^{ \lambda }\frac{x^{n- 1} }{\left( n-1 \right)!  } e^{-x}  \,\mathrm{d} x\\ 
       &= - \frac{ \lambda ^{n} }{n! }e^{- \lambda }+ \int_{0}^{ \lambda } \frac{x^{n-1} }{\left( n-1 \right)!  }e^{-x}\,\mathrm{d} x 
   \end{aligned}
     \]移项并做 有限 \(  n  \)次 归纳, 可得 \[
     \int_{0}^{ \lambda } e^{-x}\,\mathrm{d} x=  \sum _{k= 1}^{n} \frac{ \lambda ^{k} }{k! }e^{- \lambda }+  \int_{0}^{ \lambda }x^{n}\frac{e^{-x} }{n! }\,\mathrm{d} x  = F\left( n \right)-e^{- \lambda }+ \int_{0}^{ \lambda }\frac{x^{n} }{n! }e^{-x}\,\mathrm{d} x  
     \]又 \[
     \int_{0}^{ \lambda }e^{-x}\,\mathrm{d} x= 1-e^{- \lambda }
     \]带入得到 \[
     1= F\left( n \right)+ \int_{0}^{ \lambda }\frac{x^{n} }{n! }e^{-x}\,\mathrm{d} x  
     \]即 \[
     F\left( n \right)=  1-\int_{0}^{ \lambda }\frac{x^{n} }{n! }e^{-x}\,\mathrm{d} x  
     \]
\end{proof}
\begin{exercise}{}{}
设 $X \sim N(\mu, \sigma^2)$, 函数 $g(x)$ 可导, 求证: 两边都存在时有
\[
E[(X - \mu)g(X)] = \sigma^2 E[g'(X)]
\]
\end{exercise}
\begin{proof}
   若 \(  E\left[ g^{\prime} \left( X \right)  \right]   \)存在, 存在常数 \(  b  \), 使得  \[
  \left( X-\mu  \right)   g\left( X \right)=  \left( X-\mu  \right) \int_{-\infty}^{X}g^{\prime} \left( x \right) \,\mathrm{d} x + b\left( X-\mu  \right) 
    \]由于 \[
    E\left( b\left( X-\mu  \right)  \right)= bE\left( X-\mu  \right)= 0  
    \]对两边求期望, 得到 \[
    \begin{aligned}
    E\left(\left( X-\mu  \right)  g\left( X \right)  \right)&= \int  \int_{-\infty}^{X}\left( X-\mu  \right)g^{\prime} \left( x \right)\,\mathrm{d} x \,\mathrm{d} P  \\ 
     &=  \int \int\left( X-\mu  \right)\chi _{\left\{ x\le X \right\}}g^{\prime} \left( x \right)\,\mathrm{d} x\,\mathrm{d} P\\ 
      &=  \int g^{\prime} \left( x \right)   \int \left( X-\mu  \right)\chi _{\left\{ x\le X \right\}} \,\mathrm{d} P\,\mathrm{d} x
    \end{aligned}
    \] 其中 \[
   \begin{aligned}
    \int \left( X-\mu  \right) \chi _{\left\{ x\le X \right\}}\,\mathrm{d} P= \int_{x}^{\infty}\left( t-\mu  \right)p_{X}\left( t \right)\,\mathrm{d} x  & =  \int_{x}^{\infty}\left( t-\mu  \right)\frac{1 }{\sqrt{2 \pi } \sigma  }e^{-\frac{\left( t-\mu  \right)^{2}  }{2 \sigma ^{2} } }\,\mathrm{d} t   \\ 
     &=  \frac{ \sigma  }{\sqrt{2 \pi } } \int_{x}^{\infty}e^{-\frac{\left( t-\mu  \right)^{2}  }{2 \sigma ^{2} } }\,\mathrm{d}  \left( \frac{\left( t-\mu  \right)^{2}  }{2 \sigma ^{2} }  \right)  \\ 
      &= \frac{ \sigma  }{\sqrt{2 \pi } } e^{-\frac{\left( x-\mu  \right)^{2}  }{2 \sigma ^{2} } }\\ 
       &=  \sigma ^{2} p_{X}\left( x \right) 
   \end{aligned}
    \]于是 \[
    \begin{aligned}
    E\left( \left( X-\mu  \right)g^{\prime} \left( X \right)   \right)&=  \sigma ^{2}\int g^{\prime} \left( x \right) p_{X}\left( x \right)\,\mathrm{d} x \\ 
     &=  \sigma ^{2}E\left( g^{\prime} \left( x \right)  \right)  
    \end{aligned} 
    \]
\end{proof}
\begin{exercise}{}{}
设 $X$ 为取值为非负整数值的随机变量, $P(s) = E(s^X)$ 为 $X$ 的母函数.
\begin{enumerate}
    \item 求 $b(n,p)$ 的母函数.
    \item 求证: 若 $X \sim b(n,p)$, 则
    \[
    E\left(\frac{1}{1+X}\right) = \frac{1-(1-p)^{n+1}}{(n+1)p}
    \]
\end{enumerate}
\end{exercise}

\begin{proof}
    \begin{enumerate}
        \item 设\(  X\sim b\left( n,p \right)   \), 则 \[
        P\left( X= k \right)= \binom{ n }{k  }p^{k}\left( 1-p \right)^{n-k}   
        \] 我们有 \[
     \begin{aligned}
        E\left( s^{X} \right)&= \sum _{k= 0}^{\infty}s^{k}P\left( X= k \right)   \\ 
         &= \sum _{k= 0}^{n}\binom{ n }{k  } \left( ps \right)^{k}  \left( 1-p \right)^{n-k} \\ 
          &= \left( ps+ \left( 1-p \right)  \right)^{n}\\ 
           &= \left( 1+ \left( s-1 \right)p  \right)^{n}  
     \end{aligned}
        \]故 \(  b\left( n,p \right)   \) 的母函数为 \[
        P\left( s \right)= \left( 1+ \left( s-1 \right)p  \right)^{n}  
        \]
        \item    \[
        E\left( \frac{1 }{1+ X }  \right)=  \sum _{k= 0}^{n} \frac{1 }{1+ k }\binom{ n }{k  }p^{k}\left( 1-p \right)^{n-k}    
        \] \[
        \frac{1 }{1+ k } \binom{ n }{k  }= \frac{n! }{\left( k+ 1 \right) !\left( n-k \right)!  } = \frac{1 }{n+ 1 }\frac{\left( n+ 1 \right)!  }{\left( k+ 1 \right)!\left( \left( n+ 1 \right)-\left( k+ 1 \right)   \right)!   } = \frac{1 }{n+ 1 }\binom{ n+ 1 }{k+ 1  }    
        \]于是 \[
        \begin{aligned}
        E\left( \frac{1 }{1+ X }  \right)&= \sum _{k= 0}^{\infty} \frac{1 }{n+ 1 }\binom{ n+ 1 }{k+ 1  }p^{k}\left( 1-p \right)^{n-k}     \\ 
         &= \frac{1 }{p } \frac{1 }{n+ 1 }\sum _{k= 0}^{\infty} \binom{ n+ 1 }{k+ 1  }p^{k+ 1}\left( 1-p \right)^{\left( n+ 1 \right)-\left( k+ 1 \right)  } \\ 
          &= \frac{1 }{p }\frac{1 }{n+ 1 }\sum _{k= 1}^{\infty}\binom{ n+ 1 }{k  }p^{k}\left( 1-p \right)^{\left( n+ 1 \right)-k }\\ 
           &= \frac{1 }{p }\frac{1 }{n+ 1 }\left( \left( p+ 1-p \right)^{n+ 1} -\left( 1-p \right)^{n+ 1}   \right)         \\ 
            &= \frac{1-\left( 1-p \right)^{n+ 1}  }{\left( n+ 1 \right)p  } 
        \end{aligned}
        \]
    \end{enumerate}
    
\end{proof}
\begin{exercise}{}{}
设 $X$ 服从柯西分布, 即 $p_X(x) = \frac{1}{\pi(1+x^2)}$, 计算 $Y = 1/X$ 的密度函数, 观察它与 $p(x)$ 的关系.
\end{exercise}
\begin{solution}
    若 \(  t> 0  \), 则  \[
P\left( Y\le t \right)= P\left( \frac{1 }{X }  \le t\right)= P\left( X\ge \frac{1 }{t } \right)+ P\left( X< 0 \right) = 1- F_{X}\left( \frac{1 }{t }  \right)+ F_{X}\left( 0 \right)   
\]若 \(  t< 0  \), 则 \[
P\left( Y\le t \right)=  P\left( X\ge  \frac{1 }{t }  \right)  = F_{X}\left( \frac{1 }{t }  \right) 
\] 由 \[
\lim_{t \to 0^{+ }}\left( 1-F_{X}\left( \frac{1 }{t }  \right)+ F_{X}\left( 0 \right)   \right)= F_{X}\left( 0 \right)  
\]于是 \(  Y  \)的一个累积分布函数为 \[
F_{Y}\left( x \right)= \begin{cases} 1-F_{X}\left( \frac{1 }{t }  \right)+ F_{X}\left( 0 \right),& t> 0\\ 
 F_{X}\left( 0 \right),& t= 0\\ 
  F_{X}\left( -\frac{1 }{t }  \right),&t< 0     \end{cases}  
\] 通过分段求导, 并约定 \(  p_{Y}\left( 0 \right)= \frac{1 }{ \pi  }    \),  我们知道 \(  Y  \)的一个概率密度函数为 \[
p_{Y}\left(  \theta  \right)= \begin{cases} \frac{1 }{t^{2} }  p_{X}\left( \frac{1 }{t }  \right) ,& t> 0\\ 
 \frac{1 }{ \pi  },& t= 0 \\ 
 \frac{1 }{t^{2} }p_{X}\left( -\frac{1 }{t }  \right),& t< 0   \end{cases} = \begin{cases} \frac{1 }{ \pi \left( t^{2}+ 1 \right)  } ,& t> 0\\ 
   \frac{1 }{ \pi  },& t= 0\\ 
    \frac{1 }{ \pi \left( t^{2}+ 1 \right)  },& t< 0   \end{cases}=  \frac{1 }{ \pi \left( t^{2}+ 1 \right)  }=  p_{X}   \left( t \right) 
\] 故 \(  p_{Y}  = p_{X}\). 
\end{solution}

\begin{exercise}{}{}
\begin{enumerate}
    \item 对于多个事件, 叙述两两独立和相互独立的数学定义.
    \item 某人扔骰子 $N$ 次, 令 $A_{ij}$ 表示第 $i$ 次和第 $j$ 次得到同样点数. 试证明: $\{A_{ij}: 1 \le i < j \le N\}$ 两两独立但不相互独立.
\end{enumerate}
\end{exercise}
\begin{solution}
    \begin{enumerate}
        \item 对于事件 \(  A_1,A_2,\cdots ,A_{n}  \), 我们称它们是两两独立的,若对于任意的 \(  1\le i\neq j\le n  \), 都有 \[
        P\left( A_{i}A_{j} \right)= P\left( A_{i} \right)P\left( A_{j} \right)   
        \]  我们称它们是相互独立的, 若 对于任意的 \(  \left\{ A_{i_1}, A_{i_2} ,\cdots , A_{i_{k}}\right\}\subseteq \left\{ A_1,\cdots ,A_{n} \right\}  \), 都有 \[
        P\left(A_{i_1}\cdots A_{i_{k}}\right)=P\left( A_{i_1} \right)\cdots P\left( A_{i_{k}} \right)  
        \]
        \item 对于任意的 \(  1\le i< j\le N  \), 我们有  \[
        P\left( A_{ij} \right)= \frac{1 }{6 }  
        \]对于所有的\(  i_1< j_1  \), \(  i_2< j_2  \), \(  i_1\neq i_2  \)或 \(  j_1\neq j_2  \). \[
        P\left( A_{i_1j_1}A_{i_2j_2} \right)= \begin{cases} \frac{1 }{6 }\cdot \frac{1 }{6 }  ,& i_1\neq i_2, j_1\neq j_2\\ 
       \frac{1 }{6 }\cdot \frac{1 }{6 }    ,&i_1= i_2, j_1\neq j_2  \\ \frac{1 }{6 }\cdot \frac{1 }{6 },& i_1\neq i_2, j_1= j_2  
        \end{cases}   = \frac{1 }{36 }
        \] 因此总有 \[
        P\left( A_{i_1j_1}A_{i_2j_2} \right)= P\left( A_{i_1j_1} \right)P\left( A_{i_2j_2} \right)   
        \] 这表明 \(  \left\{ A_{ij} \right\}  \)是两两独立的.   \[
        \prod _{1\le i< j\le N}P\left( A_{ij} \right)= \frac{1 }{6^{\frac{1 }{2 }N\left( N-1 \right)  } }  
        \] \[
        P\left( \bigcap _{1\le i< j\le N}A_{ij} \right)= 6\cdot  \frac{1 }{6^{N} }= \frac{1 }{6^{N-1} }   
        \]二者不相等, 故它们不是相互独立的.
    \end{enumerate}
    
\end{solution}

\end{document}